% !PRINTED 636
\weberpage{636}
%<<CONT
und wenn man mit $b_0^{e-1}$ multipliziert und $b_0\eta = \omega$ setzt, so erh\"alt man
eine Gleichung von der Form (1).

5.\ \gesperrt{Eine Funktion $\omega$ des K\"orpers $\Omega$ ist eine ganze Funktion, wenn
sie irgend einer Gleichung}

\weq{5}{\psi(\omega) = \omega^m + c_1\omega^{m-1} + \cdots + c_{m-1}\omega + c_m = 0}

\gesperrt{gen\"ugt, in der $c_1, c_2, \ldots, c_m$ ganze rationale Funktionen sind, wenn
dies auch nicht die Gleichung niedrigsten Grades ist.}

Denn ist $\varphi(\omega) = 0$ die Gleichung niedrigsten Grades f\"ur $\omega$, so mu\ss{}
$\psi(t)$ f\"ur ein variables $t$ durch $\varphi(t)$ teilbar sein. Das ist einfach der
Inhalt der Gleichung (5). Es ist also

\[
\psi(t) = \chi(t)\varphi(t),
\]

worin $\psi(t)$ ebenfalls eine ganze Funktion von $t$ ist.
% [sic] printed psi(t), error for chi(t): the clause describes the factor chi(t) just introduced in psi(t) = chi(t) phi(t), and psi(t) is already integral by hypothesis
Wenn nun, wie vorausgesetzt,
$\psi(t)$ ganze Koeffizienten hat, so m\"ussen nach einem Satz von \gesperrt{Gauss} auch
die Koeffizienten von $\varphi(t)$ und von $\chi(t)$ ganze Funktionen sein. Vgl.\ Bd.\ I,
\S\,2, \S\,20, \S\,148\,$^1$).

Hieraus ergibt sich der folgende Hauptsatz:

6.\ \gesperrt{Summe, Differenz, Produkt zweier ganzer Funktionen in $\Omega$ sind wieder
ganze Funktionen.}

Sind n\"amlich $\omega'$, $\omega''$ zwei ganze Funktionen in $\Omega$, die den
Gleichungen gen\"ugen:

\weq{6}{
\begin{aligned}
\omega'^{\,n'} + b_1'\omega'^{\,n'-1} + \cdots + b_{n'-1}'\omega' + b_{n'}' &= 0, \\
\omega''^{\,n''} + b_1''\omega''^{\,n''-1} + \cdots + b_{n''-1}''\omega'' + b_{n''}'' &= 0,
\end{aligned}
}

so bezeichne man mit

\[
\omega_1, \omega_2, \ldots, \omega_m
\]

die Produkte

\[
\omega'^{\,\nu'}\omega''^{\,\nu''}; \quad \nu' = 0, 1, \ldots, n'-1, \quad
\nu'' = 0, 1, \ldots, n''-1,
\]

und mit $\omega$ eine der drei Funktionen $\omega' \pm \omega''$, $\omega'\omega''$, dann
kann man mit Hilfe der Gleichungen (6) setzen:

\weq{7}{
\begin{aligned}
\omega\omega_1 &= x_{1,1}\omega_1 + \cdots + x_{1,m}\omega_m, \\
&\cdots\cdots\cdots\cdots\cdots\cdots\cdots \\
\omega\omega_m &= x_{m,1}\omega_1 + \cdots + x_{m,m}\omega_m,
\end{aligned}
}

%CONT>>
\begin{quote}
$^1$) Bei diesem Beweis braucht nicht die Zerlegbarkeit der ganzen rationalen Funktionen in
lineare Faktoren vorausgesetzt zu werden, sondern nur die Zerlegbarkeit in irreduzible
Faktoren. Man kann also f\"ur die Zahlenkoeffizienten noch einen beliebigen
Rationalit\"atsbereich festsetzen. Wir werden aber in der Folge die Zerlegbarkeit einer
ganzen rationalen Funktion von $z$ in lineare Faktoren, d.\ h.\ den Fundamentalsatz der
Algebra, doch voraussetzen m\"ussen.
\end{quote}
